\begin{ex}
Tính kết quả của phép tính: $58 - 23$.
\shortans{35}
\loigiai{
Hàng đơn vị: $8 - 3 = 5$. Hàng chục: $5 - 2 = 3$. Vậy $58 - 23 = 35$.
\textbf{Đáp số:} $35$.
}
\end{ex}

\begin{ex}
Tính kết quả của phép tính: $76 - 34$.
\shortans{42}
\loigiai{
Hàng đơn vị: $6 - 4 = 2$. Hàng chục: $7 - 3 = 4$. Vậy $76 - 34 = 42$.
\textbf{Đáp số:} $42$.
}
\end{ex}

\begin{ex}
Tính kết quả của phép tính: $89 - 45$.
\shortans{44}
\loigiai{
Hàng đơn vị: $9 - 5 = 4$. Hàng chục: $8 - 4 = 4$. Vậy $89 - 45 = 44$.
\textbf{Đáp số:} $44$.
}
\end{ex}

\begin{ex}
Điền số thích hợp: $67 - \dots = 32$.
\shortans{35}
\loigiai{
Số cần điền là: $67 - 32 = 35$.
\textbf{Đáp số:} $35$.
}
\end{ex}

\begin{ex}
Điền số thích hợp: $\dots - 24 = 53$.
\shortans{77}
\loigiai{
Số cần điền là: $53 + 24 = 77$.
\textbf{Đáp số:} $77$.
}
\end{ex}

\begin{ex}
Trường Tiểu học có 85 bạn học sinh khối 1. Trong đó có 42 bạn nam. Hỏi khối 1 có bao nhiêu bạn nữ?
\shortans{43}
\loigiai{
Số học sinh nữ khối 1 là: $85 - 42 = 43$ (bạn).
\textbf{Đáp số:} $43$.
}
\end{ex}

\begin{ex}
Một cửa hàng có 96 chiếc quạt điện, đã bán được 51 chiếc quạt. Hỏi cửa hàng còn lại bao nhiêu chiếc quạt điện?
\shortans{45}
\loigiai{
Số quạt điện còn lại là: $96 - 51 = 45$ (chiếc quạt).
\textbf{Đáp số:} $45$.
}
\end{ex}

\begin{ex}
Bác nông dân có 78 kg gạo, bác đã bán đi 35 kg gạo. Hỏi bác nông dân còn lại bao nhiêu ki-lô-gam gạo?
\shortans{43}
\loigiai{
Số ki-lô-gam gạo còn lại là: $78 - 35 = 43$ (kg).
\textbf{Đáp số:} $43$.
}
\end{ex}

\begin{ex}
Một sợi dây thừng dài 88 cm. Người ta cắt đi một đoạn dài 36 cm. Hỏi sợi dây thừng còn lại dài bao nhiêu xăng-ti-mét?
\shortans{52}
\loigiai{
Chiều dài còn lại của sợi dây là: $88 - 36 = 52$ (cm).
\textbf{Đáp số:} $52$.
}
\end{ex}

\begin{ex}
Tính giá trị của biểu thức: $98 - 32 - 24$.
\shortans{42}
\loigiai{
Thực hiện: $98 - 32 = 66$; $66 - 24 = 42$.
\textbf{Đáp số:} $42$.
}
\end{ex}

\begin{ex}
Tính giá trị của biểu thức: $75 - 20 - 15$.
\shortans{40}
\loigiai{
Thực hiện: $75 - 20 = 55$; $55 - 15 = 40$.
\textbf{Đáp số:} $40$.
}
\end{ex}

\begin{ex}
Tìm số tự nhiên $x$ thỏa mãn: $x + 31 = 86$.
\shortans{55}
\loigiai{
Ta có: $x = 86 - 31 = 55$.
\textbf{Đáp số:} $55$.
}
\end{ex}

\begin{ex}
Tìm số tự nhiên $y$ thỏa mãn: $y + 43 = 97$.
\shortans{54}
\loigiai{
Ta có: $y = 97 - 43 = 54$.
\textbf{Đáp số:} $54$.
}
\end{ex}

\begin{ex}
Đàn vịt có 65 con. Có 32 con vịt bơi sang ao bên cạnh. Hỏi đàn vịt còn lại bao nhiêu con vịt?
\shortans{33}
\loigiai{
Số con vịt còn lại là: $65 - 32 = 33$ (con).
\textbf{Đáp số:} $33$.
}
\end{ex}

\begin{ex}
Một bãi đỗ xe có 59 chiếc ô tô, có 24 chiếc ô tô rời khỏi bãi. Hỏi trong bãi còn lại bao nhiêu chiếc ô tô?
\shortans{35}
\loigiai{
Số ô tô còn lại là: $59 - 24 = 35$ (chiếc ô tô).
\textbf{Đáp số:} $35$.
}
\end{ex}

\begin{ex}
Hiệu giữa số lớn nhất có hai chữ số và số 34 bằng bao nhiêu?
\shortans{65}
\loigiai{
Số lớn nhất có hai chữ số là 99: $99 - 34 = 65$.
\textbf{Đáp số:} $65$.
}
\end{ex}

\begin{ex}
Hiệu giữa số 87 và số tròn chục nhỏ nhất có hai chữ số là bao nhiêu?
\shortans{77}
\loigiai{
Số tròn chục nhỏ nhất có hai chữ số là 10: $87 - 10 = 77$.
\textbf{Đáp số:} $77$.
}
\end{ex}

\begin{ex}
Một cuốn sách có 95 trang. Mai đã đọc được 43 trang. Hỏi Mai còn bao nhiêu trang chưa đọc?
\shortans{52}
\loigiai{
Số trang Mai chưa đọc là: $95 - 43 = 52$ (trang).
\textbf{Đáp số:} $52$.
}
\end{ex}

\begin{ex}
Hai bạn An và Bình có tổng cộng 78 viên bi. Biết An có 36 viên bi. Hỏi Bình có bao nhiêu viên bi?
\shortans{42}
\loigiai{
Số viên bi của Bình là: $78 - 36 = 42$ (viên bi).
\textbf{Đáp số:} $42$.
}
\end{ex}

\begin{ex}
Số bi của Bình nhiều hơn số bi của An ở bài toán trên là bao nhiêu viên bi?
\shortans{6}
\loigiai{
Số bi của Bình nhiều hơn của An là: $42 - 36 = 6$ (viên bi).
\textbf{Đáp số:} $6$.
}
\end{ex}

\begin{ex}
Tìm số thích hợp: $77 - 23 = 11 + \dots$.
\shortans{43}
\loigiai{
Vế trái: $77 - 23 = 54$. Số cần điền: $54 - 11 = 43$.
\textbf{Đáp số:} $43$.
}
\end{ex}

\begin{ex}
So sánh: $86 - 41$ và $97 - 52$. Kết quả của mỗi phép tính đều bằng bao nhiêu?
\shortans{45}
\loigiai{
Ta có: $86 - 41 = 45$ và $97 - 52 = 45$. Cả hai phép tính đều bằng $45$.
\textbf{Đáp số:} $45$.
}
\end{ex}

\begin{ex}
Một rổ cam có 68 quả. Mẹ bán đi 25 quả và biếu bà 12 quả. Hỏi trong rổ còn lại bao nhiêu quả cam?
\shortans{31}
\loigiai{
Tổng số cam đã bán và biếu là: $25 + 12 = 37$ (quả).\nSố cam còn lại là: $68 - 37 = 31$ (quả cam).
\textbf{Đáp số:} $31$.
}
\end{ex}

\begin{ex}
Tính: $69 - 14 - 23$.
\shortans{32}
\loigiai{
Thực hiện: $69 - 14 = 55$; $55 - 23 = 32$.
\textbf{Đáp số:} $32$.
}
\end{ex}

\begin{ex}
Một cây tre dài 98 cm, chặt bớt đi một đoạn dài 45 cm. Đoạn tre còn lại dài bao nhiêu xăng-ti-mét?
\shortans{53}
\loigiai{
Đoạn tre còn lại dài: $98 - 45 = 53$ (cm).
\textbf{Đáp số:} $53$.
}
\end{ex}
